% UNREGISTERED D.70 candidate. Included sources remain frozen for checkpoint09.
% Requires root and independent proof audit before promotion.
\subsection{Generic cusp rigidity over a number field}

\begin{statement}{D.70}{Draft}{Generic full pro-prime cusp fibres on affine curves and Fermat covers}
Let \(k\) be a number field, let \(U/k\) be a smooth geometrically
connected affine hyperbolic curve, and suppose that its smooth
compactification \(X\) has at least three geometric boundary
points and at least one \(k\)-rational boundary point.
Fix such a rational boundary point \(d\) and a nonzero rational
tangent there. No other boundary point is required to be rational. Use its
actual tangential section \(s_0\) as base point for section cocycles.
For each prime \(\ell\), put
\[
 P_\ell=\pi_1(U_{\bar k})^{(\ell)},\qquad
 V_\ell=P_\ell^{\rm ab},\qquad
 {\cal K}_{k,\ell}=H^1(G_k,\mathbb Z_\ell(1)).
                                                        \tag{D.70.1}
\]
The peripheral inclusion at \(d\) gives an injection
\[
 \iota_{\ell,*}:{\cal K}_{k,\ell}
                    \longrightarrow H^1(G_k,V_\ell).   \tag{D.70.2}
\]
Give \({\cal K}_{k,\ell}\) its intrinsic \(\ell\)-adic topology.
There is a dense \(G_\delta\) subset
\({\cal G}_{k,\ell,d}\subset{\cal K}_{k,\ell}\) such that, for
every \(a\) in this subset, the fibre
\[
 H^1_{\rm cts}(G_k,P_\ell)\longrightarrow H^1(G_k,V_\ell)
           \quad\hbox{over }\iota_{\ell,*}(a)             \tag{D.70.3}
\]
is the singleton consisting of the actual completed cusp section
with parameter \(a\). The subset can be chosen so that every
parameter is ramified at infinitely many finite places of \(k\).
This assertion uses no local pro-\(\ell\) section conjecture.
No corresponding assertion for one or two geometric cusps is made.

Give
\[
 \widehat{\cal K}_k=H^1(G_k,\widehat{\mathbb Z}(1))
                    =\prod_\ell{\cal K}_{k,\ell}         \tag{D.70.4}
\]
the product topology. There is a dense \(G_\delta\) subset of this
space for which the corresponding \emph{full profinite} abelian
section fibre is also the singleton cusp class. Every primewise
parameter can simultaneously be infinitely ramified. This last
assertion uses D.24 and its proved local input D.56.

In particular, this applies over \(\mathbb Q\) to every odd
Fermat curve
\[
 U_n=\{X^n+Y^n+Z^n=0,\ XYZ\ne0\},\qquad n\ge3,
                                                        \tag{D.70.5}
\]
at each of its three rational cusps, and at every coefficient
prime \(\ell\), including those not dividing \(n\). For the
\(3\)-power and \(5\)-power curves of D.20 it supplies generic
singleton full pro-\(\ell\) fibres at their actual cusp abelian
classes. It does not classify the metabelian fibres in D.20
or give a first finite obstruction for their nonzero members.
\end{statement}

\begin{proof}
\paragraph{The integral peripheral module.}
Write \(D=X-U\), let \(r=|D(\bar k)|\ge3\), and fix \(\ell\).
The group \(P=P_\ell\) is free pro-\(\ell\) of rank \(2g+r-1\).
Every cusp generator is primitive: in the usual surface
presentation eliminate a different peripheral generator.
Write \(c\) for the primitive generator at \(d\), \(H=P'\),
and \(V=P^{\rm ab}\). The boundary sequence is an exact
sequence of free integral \(\mathbb Z_\ell\)-modules
with continuous Galois action,
\[
 0\longrightarrow B
   \longrightarrow V\longrightarrow T_\ell J_X
   \longrightarrow0,\qquad
 B=\left(\bigoplus_{e\in D(\bar k)}\mathbb Z_\ell e\right)
                         \Big/\mathbb Z_\ell\sum_e e\ (1).
                                                        \tag{D.70.6}
\]
Here the final \((1)\) twists the whole permutation quotient.
Let \(I=\mathbb Z_\ell(1)c\subset B\). It is primitive, and
\(B/I\) is free of rank \(r-2\); hence \(V/I\) is also
torsion-free. Any two distinct peripheral lines have zero
intersection in \(B\). Indeed, if \(a e-b f\) is a multiple
of \(\sum_e e\), a third cusp coordinate forces that multiple
to be zero, and then \(a=b=0\).

Choose a finite set \(S\) of places containing the archimedean
places, the places above \(\ell\), all bad reduction places of
the marked curve \((X,D)\), and the places at which the chosen
tangent is not integral and nonzero on the special fibre.
At \(v\notin S\), the marked model is smooth with etale
boundary. A finite unramified extension of \(k_v\) splits
all its boundary points. On that extension \(B/I\) is a
direct sum of copies of \(\mathbb Z_\ell(1)\), and its
Frobenius acts by a positive power of \(N(v)\), which is
not one. Thus \((B/I)^{G_{k_v}}=0\), integrally.
Also
\[
                  (T_\ell J_X)^{G_{k_v}}=0.             \tag{D.70.7}
\]
For this, good reduction and \(\ell\ne\operatorname{char}
\kappa(v)\) inject \(J_X(k_v)[\ell^\infty]\) into the
finite group \(J_X(\kappa(v))\): the kernel of reduction
is pro-residue-characteristic. A finite \(\ell\)-primary
group has zero Tate module. Taking invariants in (D.70.6),
and then in its quotient by \(I\), gives
\[
 V^{G_{k_v}}=(V/I)^{G_{k_v}}=0.
                                                        \tag{D.70.8}
\]
The same vanishings hold for \(G_k\), by restriction to
one such decomposition group. The continuous cohomology
sequence of \(0\to I\to V\to V/I\to0\) therefore proves
both (D.70.2) and the local injections
\[
 H^1(k_v,I)\lhook\joinrel\longrightarrow H^1(k_v,V)
                    \qquad(v\notin S).                 \tag{D.70.9}
\]
No rationalization is used: the coefficient sequence
has a continuous integral linear splitting after
forgetting the Galois action.

\paragraph{Completed multiplicative groups and actual cocycles.}
The finite Kummer sequences and the Mittag--Leffler
property of the finite groups \(H^0(k,\mu_{\ell^n})\) give
\[
 {\cal K}_{k,\ell}=\varprojlim_n
                         k^\times/(k^\times)^{\ell^n}. \tag{D.70.10}
\]
Here each finite-exponent quotient is discrete and need
not be finite. We make the topology and representatives
explicit. The principal divisor subgroup
\(\operatorname{Prin}(k)\subset\bigoplus_{v\ {\rm finite}}
\mathbb Z v\) has finite index, with quotient the ideal
class group; its finiteness is
\cite[Theorem 4.4]{DE_MilneANT}. It is free abelian.
The exact sequence
\[
 1\longrightarrow\mathcal O_k^\times
 \longrightarrow k^\times\longrightarrow
             \operatorname{Prin}(k)\longrightarrow0
                                                        \tag{D.70.11}
\]
splits as abstract abelian groups because its quotient
is free. The unit theorem
\cite[Theorem 5.1]{DE_MilneANT} then gives a decomposition
\[
       k^\times=\mu(k)\times L,\qquad
       L\simeq\bigoplus_{i\in{\cal I}}\mathbb Z,
                                                        \tag{D.70.12}
\]
where \({\cal I}\) is countable; its finite unit-rank
part is included in this index set. Fix one such
decomposition. Consequently
\[
 {\cal K}_{k,\ell}\simeq
       \mu(k)(\ell)\times c_0({\cal I},\mathbb Z_\ell).
                                                        \tag{D.70.13}
\]
The finite root-of-unity factor is discrete, and the
second factor has the uniform \(\ell\)-adic topology.
This is the topology in (D.70.10), and also the intrinsic
topology with neighborhoods \(\ell^n{\cal K}_{k,\ell}\);
it is independent of the chosen splitting. It is a
complete separable metric topology.

Fix compatible roots for every generator of \(L\),
and thereby an integral continuous Kummer cocycle
\(\kappa_i:G_k\to\mathbb Z_\ell(1)\). Their uniformly
convergent sums
\(\sum_i a_i\kappa_i\), for \((a_i)\in c_0\), give a
continuous linear choice of actual cocycles on the
free factor. For each of the finitely many torsion
classes in \(\mu(k)(\ell)\), choose one actual integral
continuous Kummer representative. On each fixed
torsion component, add that representative to the
free-factor sum. This gives a continuous choice of
actual cocycle \(a(g)\) for every parameter. We do
not assert a linear section on the torsion factor,
at two or at any other prime.

For each finite place let \(a_v\in\mathbb Z_\ell\)
be its valuation coordinate. The valuation map is
continuous, and for each \(m\) only finitely many
\(a_v\) are nonzero modulo \(\ell^m\). To see this
directly, the class modulo \(\ell^m\) in (D.70.10)
has a representative in \(k^\times\), whose divisor
has finite support. At \(v\nmid\ell\), restriction
of the actual cocycle to inertia is precisely
\(a_v\) times the integral tame \(\ell\)-character.
Roots of unity and local units have zero such inertia.

Let \(h\) be the class number of \(k\). For every
finite \(v\), choose \(\alpha_v\in k^\times\) with
\[
                       \operatorname{div}(\alpha_v)=h v.
                                                        \tag{D.70.14}
\]
Multiply it by its inverse torsion component in
(D.70.12), so that \(\alpha_v\in L\) without changing
its divisor. Its cocycle is now the prescribed finite
integral linear combination of the \(\kappa_i\).
The perturbation \(\ell^M\kappa(\alpha_v)\) preserves
the torsion component, is uniformly divisible by
\(\ell^M\) as an actual cochain, and changes only
the valuation at \(v\), by \(h\ell^M\).

\paragraph{The actual local cusp comparison.}
The rational tangential action
\(\sigma_0:G_k\to\operatorname{Aut}(P)\) satisfies
\(\sigma_0(g)(c)=c^{\chi_\ell(g)}\). Its restriction
at \(v\notin S\) is unramified. This is the unit
tangential section of the marked good-reduction
specialization; a unit tangent has trivial
prime-to-residue-characteristic Kummer inertia.
The completed parameter gives the actual action
\[
 \sigma_a(g)=\operatorname{Int}(c^{a(g)})\sigma_0(g).
                                                        \tag{D.70.15}
\]
Every section class over \(\iota_{\ell,*}(a)\) can
be conjugated to have exactly that abelian cocycle.
Its difference from this cusp section is an actual
nonabelian cocycle \(z:G_k\to H\) for \(\sigma_a\).
Conversely every such cocycle supplies a section
in the fibre. By \(V^{G_k}=0\), every conjugator
between two exact representatives lies in \(H\).
Thus the fibre in (D.70.3) is exactly
\(H^1(G_k,H;\sigma_a)\).

Fix \(v\notin S\). Use the actual cartesian
specialization square of
\cite[Definition A.3 and Remark A.3.1, p.~38]{DE_Hoshi2012}
with \(\Sigma=\{\ell\}\). Proposition A.6 of that
source says that the specialized inertia of every
actual local section is zero or a nonzero subgroup
of one cusp inertia group. This is a statement for
ordinary geometrically pro-\(\ell\) sections of
every hyperbolic curve with good reduction, as the
setup on pp.~36--37 specifies.

For completeness the further group step is the
same integral one used in D.67. A primitive cusp
generator in a free pro-\(\ell\) group has
\[
 C_P(c^b)=\langle c\rangle
                    \quad(0\ne b\in\mathbb Z_\ell).
                                                        \tag{D.70.16}
\]
Choose a free basis containing \(c\). In its Magnus
embedding, commuting with \((1+C)^b\) implies
commuting with \(C\), because
\((1+C)^b-1\) has a compositional inverse over
\(\mathbb Q_\ell\). Comparing coefficients of
words involving the other basis variables gives
the centralizer \(\mathbb Q_\ell[[C]]\) in the
completed associative algebra. Retraction killing
those other variables then identifies a group
element in it as the power of \(c\) specified by
its integral first abelian coordinate. A geometric
normalizer of a nonzero subgroup of \(\langle c\rangle\)
centralizes it: its exponent automorphism is
forced to be one by abelianization.

If \(a_v\ne0\), the actual section has nonzero
specialized inertia, and its abelian image has
line \(\mathbb Z_\ell c\). The disjointness of
the geometric peripheral lines proved after
(D.70.6) forces the cusp in Hoshi's alternative
to be precisely \(d\). This comparison can also
be made after the finite unramified extension
splitting all boundary points; inertia is
unchanged by that extension. It distinguishes
every geometric cusp, including the extra
nonrational ones.

The whole specialized section normalizes its
inertia image. Comparing it with the rational
tangential lift at \(d\), (D.70.16) puts its
image in that cusp decomposition group.
The cartesian square pulls this conclusion
back to the actual local cusp packet. By
(D.70.9), its parameter equals \(a\) in
\(H^1(k_v,I)\). First conjugate into the
packet, and then conjugate by a cusp-inertia
element to adjust its Kummer representative
to the reference one. The resulting
conjugator has \(G_{k_v}\)-fixed abelian
image since the two actual abelian cocycles
agree. Equation (D.70.8) makes that image
zero, so the conjugator lies in \(H\).
If \(a_v=0\), Hoshi's nonzero alternative
would have nonzero abelian inertia, which
is impossible; both sections have zero
specialized inertia. We have proved the
two exact assertions
\[
 \begin{cases}
 z|_{G_{k_v}}\text{ is an }H\text{-coboundary},
                              &a_v\ne0,\\
 z(I_v)=1,                    &a_v=0,
 \end{cases}\qquad v\notin S.                         \tag{D.70.17}
\]
This argument uses the integral quotient \(V/I\),
not a rationalized abelian Kummer map.

\paragraph{Finite capture and fresh-prime witnesses.}
Let \(P_0=P\), \(P_{m+1}=\Phi(P_m)\), and put
\[
 Q_m=P/P_m,\quad H_m=\operatorname{im}(H\to Q_m),\quad
 S_m(a)=S\cup\{v:a_v\notin\ell^m\mathbb Z_\ell\}.
                                                        \tag{D.70.18}
\]
The \(P_m\) are characteristic open subgroups
with trivial intersection, and \(Q_m\) has
exponent dividing \(\ell^m\). These finite
quotients induce the topology of \(H\).
Write \(\rho_{a,m}\) for its actual action on
\(H_m\). Both \(S_m(a)\) and this action are
constant when the torsion component of \(a\)
is fixed and its free component is fixed
modulo \(\ell^m\).

Every actual cocycle \(z:G_k\to H\) reduces
to the finite set
\[
 E_m(a)=Z^1(G_{k,S_m(a)},H_m;\rho_{a,m}).
                                                        \tag{D.70.19}
\]
Indeed outside \(S_m(a)\) its inertia is zero
if \(a_v=0\). Otherwise (D.70.17) makes it
an integral conjugation coboundary; the
inertia action is
\(\operatorname{Int}(c^{a(\tau)})\), which is
trivial on \(Q_m\) because \(a(\tau)\) is
divisible by \(\ell^m\). The finite action
itself has this same support bound. Thus
the reduced \emph{actual} affine
representation factors through the stated
finite-ramification Galois group.
The set in (D.70.19) is finite by the
finite-extension finiteness theorem: its
cocycles are homomorphisms into the fixed
finite group \(H_m\rtimes\operatorname{im}\rho_{a,m}\),
with fixed second coordinate and ramification
in the fixed finite set \(S_m(a)\).

Call \(z\in E_m(a)\) exceptional if its
restriction to \(\ker\rho_{a,m}\) is nontrivial.
Let \({\cal V}_m\) be the set of parameters
for which each exceptional \(z\) has a
witness place \(v\notin S_m(a)\) with
\[
 \rho_{a,m}(\operatorname{Frob}_v)=1,\qquad
 z(\operatorname{Frob}_v)\ne1,\qquad a_v\ne0.
                                                        \tag{D.70.20}
\]
This is open: all finite data are locally
constant, there are only finitely many
exceptional cocycles, and nonvanishing of
finitely many specified valuation coordinates
is open.

It is dense on every fixed torsion component.
Given a neighborhood with free coordinates
fixed modulo \(\ell^N\), choose
\(M\ge\max\{m,N\}\) and write \(w=v_\ell(h)\).
For an exceptional \(z\), a group element
\(\delta\in\ker\rho_{a,m}\) with
\(z(\delta)\ne1\) gives the actual element
\((z(\delta),1)\) in its finite affine image.
Chebotarev, in precisely the finite-image
form of \cite[Chapter V, Theorem 3.23]{PBC_MilneCFT},
provides infinitely many places in its
conjugacy class. The two finite tests in
(D.70.20) are invariant under conjugacy.
Choose distinct such places for the
finitely many exceptional cocycles, avoiding
\(S_m(a)\) and the finite set where
\(a_v\notin\ell^{M+w+1}\mathbb Z_\ell\).
Add \(\ell^M\kappa(\alpha_v)\) at each
chosen place. By (D.70.14), the new
coordinate at that place has valuation
exactly \(M+w\); in particular it is
nonzero. All other chosen valuation
coordinates remain unchanged. The actual
finite action, support set, cocycle list,
torsion component, and prescribed
neighborhood are preserved. This proves
density. No principal-prime restriction
has been imposed on Chebotarev.

For every finite set \(T\) of finite places,
the locus \(a_v=0\) for all \(v\notin T\)
is closed and has empty interior: a fresh
perturbation (D.70.14), at sufficiently
large depth, leaves any neighborhood.
There are countably many such \(T\).
The Baire theorem applied to (D.70.13)
therefore gives a dense \(G_\delta\) set
\[
 {\cal G}_{k,\ell,d}\subset\bigcap_{m\ge1}{\cal V}_m
                                                        \tag{D.70.21}
\]
whose parameters all have infinite valuation
support. In particular each has \(a_v\ne0\)
at some \(v\notin S\).

\paragraph{Descent and the normal inner inertia subgroup.}
Fix \(a\) in (D.70.21) and an actual
cocycle \(z:G_k\to H\). None of its
reductions in (D.70.19) is exceptional.
Otherwise its witness in (D.70.20) and
the local coboundary assertion (D.70.17)
would give simultaneously a nontrivial
and a trivial value at a Frobenius with
identity action. Hence \(z\) is trivial
on the kernel of the actual action on
\(H\), and descends continuously to
\[
 {\cal I}_a=\operatorname{im}\left(
     G_k\longrightarrow\prod_m\operatorname{Aut}(H_m)\right).
                                                        \tag{D.70.22}
\]
This is the compact actual action image.
Continuity follows from the quotient map
of compact groups.

Choose \(v\notin S\) with \(a_v\ne0\).
Its inertia image in (D.70.22) is
\[
 N_v=\{\operatorname{Int}(c^{a_v t})|_H:
                                  t\in\mathbb Z_\ell\}.
                                                        \tag{D.70.23}
\]
It is normal in \({\cal I}_a\), since every
operator in (D.70.15) sends \(c\) to
\(c^{\chi_\ell(g)}\), and \(\chi_\ell(g)\)
is a unit. Equation (D.70.16) and the
primitive abelian image of \(c\) imply
\[
                    H^{N_v}=H\cap\langle c\rangle=1.
                                                        \tag{D.70.24}
\]
This includes nonprimitive \(a_v\); no
division by \(a_v\) is used in \(H\).
By (D.70.17) the descended cocycle on
\(N_v\) is an integral \(H\)-coboundary.
Conjugate it by that element of \(H\) to
make it trivial on \(N_v\). Comparing its
values at \(ng=g(g^{-1}ng)\) shows that
every value is fixed by \(N_v\), hence
is one by (D.70.24). Thus the whole
cocycle is a coboundary. This proves
the asserted singleton in (D.70.3).

\paragraph{The full profinite fibre and the actual Fermat maps.}
Continuous cochains into the product
\(\widehat{\mathbb Z}(1)=\prod_\ell\mathbb Z_\ell(1)\)
are exactly products of continuous
component cochains, and the same is true
of coboundaries. This proves (D.70.4).
It is also the power-completion topology
\(\varprojlim_n k^\times/(k^\times)^n\)
with discrete quotients. The countable
product of the complete separable metric
spaces in (D.70.13) is again such a space.
The product of the sets (D.70.21) is a
dense \(G_\delta\): write each as a
countable intersection of open dense
sets and pull those sets back by the
coordinate projections.

For every parameter in this product,
the product of the chosen component
Kummer cocycles gives an actual continuous
\(\widehat{\mathbb Z}(1)\)-valued cocycle,
hence a full cusp section. Any other
actual full section with its abelian
class has, by (D.70.3), the same class
as this one in every geometric pro-\(\ell\)
quotient. Apply precisely D.24, with its
D.56 local input, to these two actual
full sections. They have the same full
class. No assertion that primewise
geometric quotients detect arbitrary
geometric elements is used.

For (D.70.5), the actual finite etale
map to the tripod is
\[
 t=-\left(\frac XZ\right)^n,\qquad
 1-t=-\left(\frac YZ\right)^n.                          \tag{D.70.25}
\]
It has degree \(n^2\); its geometric
deck group is \(\mu_n^2\). The three
rational cusps are \([0:-1:1]\),
\([-1:0:1]\), and \([1:-1:0]\).
There are \(3n\) geometric cusps, so
the preceding local argument distinguishes
also all the additional nonrational ones.
At \(O=[0:-1:1]\), with coordinate
\(u=X/Z\), the local map is exactly
\(t=-u^n\). Consequently its actual
cusp abelian parameter maps to the
tripod pair
\[
                  \bigl(\kappa(-1)+na,\ 0\bigr).
                                                        \tag{D.70.26}
\]
The sign term is retained at two.
This identifies the cusp and its
abelian map with the ones used in
D.20, rather than introducing an
abstract free-group model. The proof
of (D.70.3) takes place in the actual
group of \(U_n\) itself, so it also
applies when \(\ell\nmid n\), where
one cannot infer injectivity merely
from the map to the tripod.
The zero-abelian metabelian families
of D.20 lie at \(a=0\), outside our
infinitely ramified generic set.
No claim about their metabelian
fibres is changed.
\end{proof}
